
\documentclass{article}
\usepackage{colortbl}
\usepackage{makecell}
\usepackage{multirow}
\usepackage{supertabular}

\begin{document}

\newcounter{utterance}

\centering \large Interaction Transcript for game `wordle', experiment `medium\_frequency\_words\_no\_clue\_no\_critic', episode 2 with deepseek{-}v4{-}pro.
\vspace{24pt}

{ \footnotesize  \setcounter{utterance}{1}
\setlength{\tabcolsep}{0pt}
\begin{supertabular}{c@{$\;$}|p{.15\linewidth}@{}p{.15\linewidth}p{.15\linewidth}p{.15\linewidth}p{.15\linewidth}p{.15\linewidth}}
    \# & \multicolumn{2}{c}{Player} && \multicolumn{2}{c}{Game Master} \\
    \hline

    \theutterance \stepcounter{utterance}  
    & & & \multicolumn{4}{p{0.6\linewidth}}{
        \cellcolor[rgb]{0.9,0.9,0.9}{
            \makecell[{{p{\linewidth}}}]{
                \texttt{\tiny{[P1$\langle$GM]}}
                \texttt{Du är en språkmagiker som gillar att gissa ord genom att använda de givna reglerna.} \\
\\ 
\texttt{Välkommen till Wordle! Du har sex försök på dig att gissa målordet, ett giltigt svenskt ord med fem gemena bokstäver. Använd taggarna "förklaring:" och "gissning:" för att ge en kortfattad förklaring till varje gissning.} \\
\\ 
\texttt{Till exempel, om din gissning är "fråga", ska ditt svar vara} \\
\texttt{förklaring: detta är ett vanligt svenskt ord med fem bokstäver, och jag börjar min gissning med detta ord.} \\
\texttt{gissning: fråga} \\
\\ 
\texttt{Efter varje gissning kommer ditt svar att valideras, och du kommer att få återkoppling som visar vilka bokstäver som är korrekta (gröna), vilka bokstäver som är korrekta men på fel position (gula), och vilka bokstäver som är felaktiga (röda). Denna återkoppling kan vara användbar för att avgöra vilka bokstäver du ska ta med eller utesluta i din nästa gissning.} \\
\\ 
\texttt{Till exempel kan återkopplingen för "fråga" vara:} \\
\texttt{gissning\_återkoppling: f<gul> r<gul> å<grön> g<gul> a<röd>} \\
\\ 
\texttt{Förklaringen ska innehålla detaljer om hur gissning\_återkoppling används för att komma fram till en ny gissning.} \\
\\ 
\texttt{Generera ingen annan text.} \\
\\ 
\texttt{Låt oss börja med din första gissning.} \\
            }
        }
    }
     \\ \\

    \theutterance \stepcounter{utterance}  
    & \multicolumn{4}{p{0.6\linewidth}}{
        \cellcolor[rgb]{0.561,0.737,0.561}{
            \makecell[{{p{\linewidth}}}]{
                \texttt{\tiny{[P1$\rangle$GM]}}
                \texttt{förklaring: Jag börjar med ett vanligt svenskt ord som innehåller många olika vokaler och konsonanter för att få en bred första återkoppling.} \\
\texttt{gissning: spela} \\
            }
        }
    }
    & & \\ \\

    \theutterance \stepcounter{utterance}  
    & & & \multicolumn{2}{p{0.3\linewidth}}{
        \cellcolor[rgb]{0.9,0.9,0.9}{
            \makecell[{{p{\linewidth}}}]{
                \texttt{\tiny{[GM$|$GM]}}
                \texttt{attempts = 1} \\
\texttt{target\_word = hylla} \\
\texttt{guess = spela} \\
\texttt{guess\_feedback = s<red> p<red> e<red> l<green> a<green>} \\
            }
        }
    }
    & & \\ \\

    \theutterance \stepcounter{utterance}  
    & & & \multicolumn{4}{p{0.6\linewidth}}{
        \cellcolor[rgb]{0.9,0.9,0.9}{
            \makecell[{{p{\linewidth}}}]{
                \texttt{\tiny{[P1$\langle$GM]}}
                \texttt{gissning\_återkoppling: s<röd> p<röd> e<röd> l<grön> a<grön>} \\
\\ 
\texttt{Ge ditt svar endast i detta format.} \\
\texttt{förklaring: detaljer} \\
\texttt{gissning: ord} \\
            }
        }
    }
     \\ \\

    \theutterance \stepcounter{utterance}  
    & \multicolumn{4}{p{0.6\linewidth}}{
        \cellcolor[rgb]{0.561,0.737,0.561}{
            \makecell[{{p{\linewidth}}}]{
                \texttt{\tiny{[P1$\rangle$GM]}}
                \texttt{förklaring: 'l' och 'a' är gröna och sitter på position 4 och 5. 's', 'p', 'e' är helt uteslutna. Jag testar ett nytt ord med 'l' och 'a' kvar på sina positioner, och introducerar vanliga konsonanter som 'r', 't', 'n' för att hitta resten.} \\
\texttt{gissning: rulta} \\
            }
        }
    }
    & & \\ \\

    \theutterance \stepcounter{utterance}  
    & & & \multicolumn{2}{p{0.3\linewidth}}{
        \cellcolor[rgb]{0.9,0.9,0.9}{
            \makecell[{{p{\linewidth}}}]{
                \texttt{\tiny{[GM$|$GM]}}
                \texttt{Error: The guessed word is not a valid word for this game.} \\
            }
        }
    }
    & & \\ \\

    \theutterance \stepcounter{utterance}  
    & & & \multicolumn{4}{p{0.6\linewidth}}{
        \cellcolor[rgb]{0.9,0.9,0.9}{
            \makecell[{{p{\linewidth}}}]{
                \texttt{\tiny{[P1$\langle$GM]}}
                \texttt{Din gissning är inte ett giltigt ord för det här spelet. Vänligen försök igen.} \\
\\ 
\texttt{Ge ditt svar endast i detta format.} \\
\texttt{förklaring: detaljer} \\
\texttt{gissning: ord} \\
            }
        }
    }
     \\ \\

    \theutterance \stepcounter{utterance}  
    & \multicolumn{4}{p{0.6\linewidth}}{
        \cellcolor[rgb]{0.561,0.737,0.561}{
            \makecell[{{p{\linewidth}}}]{
                \texttt{\tiny{[P1$\rangle$GM]}}
                \texttt{förklaring: 'l' och 'a' är gröna på position 4 och 5. 's', 'p', 'e' är uteslutna. Jag försöker med ett vanligt ord som slutar på 'la', innehåller en ny vokal 'o' och konsonanterna 'k' och 'n'.} \\
\texttt{gissning: knyta} \\
            }
        }
    }
    & & \\ \\

    \theutterance \stepcounter{utterance}  
    & & & \multicolumn{2}{p{0.3\linewidth}}{
        \cellcolor[rgb]{0.9,0.9,0.9}{
            \makecell[{{p{\linewidth}}}]{
                \texttt{\tiny{[GM$|$GM]}}
                \texttt{attempts = 2} \\
\texttt{target\_word = hylla} \\
\texttt{guess = knyta} \\
\texttt{guess\_feedback = k<red> n<red> y<yellow> t<red> a<green>} \\
            }
        }
    }
    & & \\ \\

    \theutterance \stepcounter{utterance}  
    & & & \multicolumn{4}{p{0.6\linewidth}}{
        \cellcolor[rgb]{0.9,0.9,0.9}{
            \makecell[{{p{\linewidth}}}]{
                \texttt{\tiny{[P1$\langle$GM]}}
                \texttt{gissning\_återkoppling: k<röd> n<röd> y<gul> t<röd> a<grön>} \\
\\ 
\texttt{Ge ditt svar endast i detta format.} \\
\texttt{förklaring: detaljer} \\
\texttt{gissning: ord} \\
            }
        }
    }
     \\ \\

    \theutterance \stepcounter{utterance}  
    & \multicolumn{4}{p{0.6\linewidth}}{
        \cellcolor[rgb]{0.561,0.737,0.561}{
            \makecell[{{p{\linewidth}}}]{
                \texttt{\tiny{[P1$\rangle$GM]}}
                \texttt{förklaring: 'a' är grön på position 5, och 'l' är grön på position 4 enligt tidigare återkoppling. 'y' är gul, så den finns i ordet men inte på position 3. 'k', 'n', 't' är uteslutna. Jag letar efter ett ord med mönstret \_ \_ \_ l a, där 'y' finns med på någon av de tre första positionerna. Ett möjligt ord är 'yngla' som betyder att föda ungar (om fiskar).} \\
\texttt{gissning: yngla} \\
            }
        }
    }
    & & \\ \\

    \theutterance \stepcounter{utterance}  
    & & & \multicolumn{2}{p{0.3\linewidth}}{
        \cellcolor[rgb]{0.9,0.9,0.9}{
            \makecell[{{p{\linewidth}}}]{
                \texttt{\tiny{[GM$|$GM]}}
                \texttt{Error: The guessed word is not a valid word for this game.} \\
            }
        }
    }
    & & \\ \\

    \theutterance \stepcounter{utterance}  
    & & & \multicolumn{4}{p{0.6\linewidth}}{
        \cellcolor[rgb]{0.9,0.9,0.9}{
            \makecell[{{p{\linewidth}}}]{
                \texttt{\tiny{[P1$\langle$GM]}}
                \texttt{Din gissning är inte ett giltigt ord för det här spelet. Vänligen försök igen.} \\
\\ 
\texttt{Ge ditt svar endast i detta format.} \\
\texttt{förklaring: detaljer} \\
\texttt{gissning: ord} \\
            }
        }
    }
     \\ \\

    \theutterance \stepcounter{utterance}  
    & \multicolumn{4}{p{0.6\linewidth}}{
        \cellcolor[rgb]{0.561,0.737,0.561}{
            \makecell[{{p{\linewidth}}}]{
                \texttt{\tiny{[P1$\rangle$GM]}}
                \texttt{förklaring: Mönstret är \_ \_ \_ l a med 'y' någonstans bland de tre första. 'yngla' var ogiltigt. Jag testar 'mylta' – ett ord som betyder att mylla jord, med 'y' på andra positionen.} \\
\texttt{gissning: mylta} \\
            }
        }
    }
    & & \\ \\

    \theutterance \stepcounter{utterance}  
    & & & \multicolumn{2}{p{0.3\linewidth}}{
        \cellcolor[rgb]{0.9,0.9,0.9}{
            \makecell[{{p{\linewidth}}}]{
                \texttt{\tiny{[GM$|$GM]}}
                \texttt{Error: The guessed word is not a valid word for this game.} \\
            }
        }
    }
    & & \\ \\

    \theutterance \stepcounter{utterance}  
    & & & \multicolumn{4}{p{0.6\linewidth}}{
        \cellcolor[rgb]{0.9,0.9,0.9}{
            \makecell[{{p{\linewidth}}}]{
                \texttt{\tiny{[P1$\langle$GM]}}
                \texttt{Din gissning är inte ett giltigt ord för det här spelet. Vänligen försök igen.} \\
\\ 
\texttt{Ge ditt svar endast i detta format.} \\
\texttt{förklaring: detaljer} \\
\texttt{gissning: ord} \\
            }
        }
    }
     \\ \\

    \theutterance \stepcounter{utterance}  
    & \multicolumn{4}{p{0.6\linewidth}}{
        \cellcolor[rgb]{0.561,0.737,0.561}{
            \makecell[{{p{\linewidth}}}]{
                \texttt{\tiny{[P1$\rangle$GM]}}
                \texttt{förklaring: Mönstret är \_ \_ \_ l a med 'y' på position 1, 2 eller 3. Bokstäverna 's', 'p', 'e', 'k', 'n', 't' är uteslutna. Jag testar 'vyala' – ett mindre vanligt ord som betyder att yla svagt, med 'y' på position 2.} \\
\texttt{gissning: vyala} \\
            }
        }
    }
    & & \\ \\

    \theutterance \stepcounter{utterance}  
    & & & \multicolumn{2}{p{0.3\linewidth}}{
        \cellcolor[rgb]{0.9,0.9,0.9}{
            \makecell[{{p{\linewidth}}}]{
                \texttt{\tiny{[GM$|$GM]}}
                \texttt{Error: The guessed word is not a valid word for this game.} \\
            }
        }
    }
    & & \\ \\

    \theutterance \stepcounter{utterance}  
    & & & \multicolumn{4}{p{0.6\linewidth}}{
        \cellcolor[rgb]{0.9,0.9,0.9}{
            \makecell[{{p{\linewidth}}}]{
                \texttt{\tiny{[P1$\langle$GM]}}
                \texttt{Din gissning är inte ett giltigt ord för det här spelet. Vänligen försök igen.} \\
\\ 
\texttt{Ge ditt svar endast i detta format.} \\
\texttt{förklaring: detaljer} \\
\texttt{gissning: ord} \\
            }
        }
    }
     \\ \\

    \theutterance \stepcounter{utterance}  
    & \multicolumn{4}{p{0.6\linewidth}}{
        \cellcolor[rgb]{0.561,0.737,0.561}{
            \makecell[{{p{\linewidth}}}]{
                \texttt{\tiny{[P1$\rangle$GM]}}
                \texttt{förklaring: Mönstret är \_ \_ \_ l a med 'y' på en av de tre första positionerna. Uteslutna bokstäver: s, p, e, k, n, t. Möjliga ord är få. Jag testar 'dylga' – ett ord som betyder gyttja eller sörja, med 'y' på position 2.} \\
\texttt{gissning: dylga} \\
            }
        }
    }
    & & \\ \\

    \theutterance \stepcounter{utterance}  
    & & & \multicolumn{2}{p{0.3\linewidth}}{
        \cellcolor[rgb]{0.9,0.9,0.9}{
            \makecell[{{p{\linewidth}}}]{
                \texttt{\tiny{[GM$|$GM]}}
                \texttt{Error: The guessed word is not a valid word for this game.} \\
            }
        }
    }
    & & \\ \\

    \theutterance \stepcounter{utterance}  
    & & & \multicolumn{4}{p{0.6\linewidth}}{
        \cellcolor[rgb]{0.9,0.9,0.9}{
            \makecell[{{p{\linewidth}}}]{
                \texttt{\tiny{[P1$\langle$GM]}}
                \texttt{Din gissning är inte ett giltigt ord för det här spelet. Vänligen försök igen.} \\
\\ 
\texttt{Ge ditt svar endast i detta format.} \\
\texttt{förklaring: detaljer} \\
\texttt{gissning: ord} \\
            }
        }
    }
     \\ \\

    \theutterance \stepcounter{utterance}  
    & \multicolumn{4}{p{0.6\linewidth}}{
        \cellcolor[rgb]{0.561,0.737,0.561}{
            \makecell[{{p{\linewidth}}}]{
                \texttt{\tiny{[P1$\rangle$GM]}}
                \texttt{förklaring: Mönstret är \_ \_ \_ l a med 'y' på en av de tre första positionerna. Uteslutna: s, p, e, k, n, t, d, g. Jag testar 'vylga' – men det är troligen ogiltigt. Bättre alternativ: 'bylta' (betyder packe eller bylte), med 'y' på position 2.} \\
\texttt{gissning: bylta} \\
            }
        }
    }
    & & \\ \\

    \theutterance \stepcounter{utterance}  
    & & & \multicolumn{2}{p{0.3\linewidth}}{
        \cellcolor[rgb]{0.9,0.9,0.9}{
            \makecell[{{p{\linewidth}}}]{
                \texttt{\tiny{[GM$|$GM]}}
                \texttt{Error: The guessed word is not a valid word for this game.} \\
            }
        }
    }
    & & \\ \\

    \theutterance \stepcounter{utterance}  
    & & & \multicolumn{4}{p{0.6\linewidth}}{
        \cellcolor[rgb]{0.9,0.9,0.9}{
            \makecell[{{p{\linewidth}}}]{
                \texttt{\tiny{[P1$\langle$GM]}}
                \texttt{Din gissning är inte ett giltigt ord för det här spelet. Vänligen försök igen.} \\
\\ 
\texttt{Ge ditt svar endast i detta format.} \\
\texttt{förklaring: detaljer} \\
\texttt{gissning: ord} \\
            }
        }
    }
     \\ \\

    \theutterance \stepcounter{utterance}  
    & \multicolumn{4}{p{0.6\linewidth}}{
        \cellcolor[rgb]{0.561,0.737,0.561}{
            \makecell[{{p{\linewidth}}}]{
                \texttt{\tiny{[P1$\rangle$GM]}}
                \texttt{förklaring: Mönstret är \_ \_ \_ l a med 'y' på position 1, 2 eller 3. Uteslutna: s, p, e, k, n, t, d, g, b. Jag testar 'fylga' – ett dialektalt ord för följe, alternativt 'mylja' men 'mylta' var ogiltigt så 'm' är osäkert. Ett bättre ord: 'yra' är för kort. 'yngla' ogiltigt. Jag testar 'ylva' – men det är 4 bokstäver. Kanske 'yvala'? Ogiltigt. Jag testar 'yppla' – men 'p' är uteslutet. Jag provar 'rulla' – men 'y' måste med. 'ylla' är 4 bokstäver. 'hylla' har inte 'y' på rätt sätt? 'hylla' har y på position 2, slutar på lla. Testar 'hylla'.} \\
\texttt{gissning: hylla} \\
            }
        }
    }
    & & \\ \\

    \theutterance \stepcounter{utterance}  
    & & & \multicolumn{2}{p{0.3\linewidth}}{
        \cellcolor[rgb]{0.9,0.9,0.9}{
            \makecell[{{p{\linewidth}}}]{
                \texttt{\tiny{[GM$|$GM]}}
                \texttt{attempts = 3} \\
\texttt{target\_word = hylla} \\
\texttt{guess = hylla} \\
\texttt{guess\_feedback = h<green> y<green> l<green> l<green> a<green>} \\
            }
        }
    }
    & & \\ \\

    \theutterance \stepcounter{utterance}  
    & & & \multicolumn{2}{p{0.3\linewidth}}{
        \cellcolor[rgb]{0.9,0.9,0.9}{
            \makecell[{{p{\linewidth}}}]{
                \texttt{\tiny{[GM$|$GM]}}
                \texttt{game\_result = WIN} \\
            }
        }
    }
    & & \\ \\

\end{supertabular}
}

\end{document}
